\makeatletter
\def\rap@fontpath{../../../Common/}
\makeatother
\documentclass{../../../Common/Latex/Templates/rapport}
% ---------------------------------------------------------------------------
% Shared settings for the H2026 General Assembly document set.
%
% Usage, directly after \documentclass{../../../Common/Latex/Templates/rapport}
% in a document one folder below this file:
%   % ---------------------------------------------------------------------------
% Shared settings for the H2026 General Assembly document set.
%
% Usage, directly after \documentclass{../../../Common/Latex/Templates/rapport}
% in a document one folder below this file:
%   % ---------------------------------------------------------------------------
% Shared settings for the H2026 General Assembly document set.
%
% Usage, directly after \documentclass{../../../Common/Latex/Templates/rapport}
% in a document one folder below this file:
%   \input{../ga-common}
%
% Image paths are relative to the compiled document's folder, so every
% document using this file must sit one folder below it.
% ---------------------------------------------------------------------------

% Meeting details, used on covers and in each document's ID block.
\newcommand{\meetingdate}{30 September 2026}
\newcommand{\meetingtime}{12:15}
\newcommand{\meetinglocation}{Festsalen, NMBU}
\newcommand{\meetingline}{\textbf{Meeting:} \meetingdate, \meetingtime, \meetinglocation}

% Header: moon logo on the left, full-text logo on the right. \projectname is
% still set because the cover title uses it.
\projectname{Eclipse NMBU}
\projectlogo{../Common/Eclipse Images/EclipseMoon}
\orglogo{../Common/Eclipse Images/EclipseFullText}

%
% Image paths are relative to the compiled document's folder, so every
% document using this file must sit one folder below it.
% ---------------------------------------------------------------------------

% Meeting details, used on covers and in each document's ID block.
\newcommand{\meetingdate}{30 September 2026}
\newcommand{\meetingtime}{12:15}
\newcommand{\meetinglocation}{Festsalen, NMBU}
\newcommand{\meetingline}{\textbf{Meeting:} \meetingdate, \meetingtime, \meetinglocation}

% Header: moon logo on the left, full-text logo on the right. \projectname is
% still set because the cover title uses it.
\projectname{Eclipse NMBU}
\projectlogo{../Common/Eclipse Images/EclipseMoon}
\orglogo{../Common/Eclipse Images/EclipseFullText}

%
% Image paths are relative to the compiled document's folder, so every
% document using this file must sit one folder below it.
% ---------------------------------------------------------------------------

% Meeting details, used on covers and in each document's ID block.
\newcommand{\meetingdate}{30 September 2026}
\newcommand{\meetingtime}{12:15}
\newcommand{\meetinglocation}{Festsalen, NMBU}
\newcommand{\meetingline}{\textbf{Meeting:} \meetingdate, \meetingtime, \meetinglocation}

% Header: moon logo on the left, full-text logo on the right. \projectname is
% still set because the cover title uses it.
\projectname{Eclipse NMBU}
\projectlogo{../Common/Eclipse Images/EclipseMoon}
\orglogo{../Common/Eclipse Images/EclipseFullText}

\projectlogo{../../../Common/Eclipse Images/EclipseMoon}
\orglogo{../../../Common/Eclipse Images/EclipseFullText}
\usepackage{enumitem}
\usepackage{array}
\usepackage{tabularx}
\usepackage{lineno}
\usepackage{hyperref}
\hypersetup{colorlinks=true, linkcolor=blue}

\renewcommand{\contentsname}{Innhold}

\begin{document}

\reportmainmatter
\begingroup
\small
\setlength{\parindent}{0pt}
\setlength{\tabcolsep}{6pt}
\renewcommand{\arraystretch}{1.2}

\begin{center}
  {\Large\bfseries Eclipse Vedtekter}\par
  \vspace{0.15cm}
  {\normalsize 26H-ECL-Vedtekter \textperiodcentered{} \textperiodcentered{} Vedtatt av: Generalforsamlingen}\par
\end{center}

\newpage
\begin{center}
  {\Large\bfseries Eclipse Vedtekter}\par
\end{center}
\tableofcontents

\linenumbers

\vsection{Del I: Grunnlag (Part I: Foundations)}

\vsubsection{\S 1 Navn og stiftelse (Name and founding)}\label{bylaw:1}
\textbf{\S 1.1 Navn (Name).}\label{bylaw:1.1} Foreningens navn er Eclipse NMBU.

\textbf{\S 1.2 Stiftelse og tilknytning (Founding and affiliation).}\label{bylaw:1.2} Eclipse NMBU er en frivillig studentforening ved Norges miljø- og biovitenskapelige universitet (NMBU), stiftet 9. oktober 2024.

\vsubsection{\S 2 Formål og grunnprinsipper (Purpose and core principles)}\label{bylaw:2}
\textbf{\S 2.1 Formål (Purpose).}\label{bylaw:2.1} Eclipse NMBU skal være en sosial og faglig arena for studenter med interesse for rommet. Foreningen er primært for studenter som er semesterregistrert ved NMBU.

\textbf{\S 2.2 Virksomhet (Activity).}\label{bylaw:2.2} Foreningens virksomhet er frivillig og rettet mot forskning og utvikling (FoU). Den overordnede visjonen er å utvikle bærekraftig romteknologi med utgangspunkt i FNs bærekraftsmål.

\textbf{\S 2.3 Grunnprinsipper (Core principles).}\label{bylaw:2.3} Foreningens virksomhet bygger på følgende grunnprinsipper:
\begin{enumerate}
  \item \textbf{Dokumentasjon:} Det som ikke er dokumentert skriftlig, anses ikke som gjort.
  \item \textbf{Ekstremt eierskap (extreme ownership):} Dersom noe går galt, ligger ansvaret både hos ledelsen og hos det enkelte medlemmet.
\end{enumerate}

\textbf{\S 2.4 Etikk og miljø (Ethics and environment).}\label{bylaw:2.4} Foreningen skal så langt det er mulig ta hensyn til miljømessige og etiske faktorer i planleggingen av sine aktiviteter. I samhandling med andre organisasjoner, foreninger, arrangementer og oppdragsgivere skal foreningen søke alternativer som ivaretar etiske og miljømessige hensyn.

\textbf{\S 2.5 Organisatorisk bærekraft (Organisational sustainability).}\label{bylaw:2.5} Foreningens struktur, oppbygging og organisering skal sikre at fremtidige generasjoner av medlemmer kan bygge videre på tidligere virksomhet, blant annet gjennom dokumentasjon, arkivering og erfaringsoverføring, jf. \S\S10.2 og 16.5.

\vsubsection{\S 3 Organisasjonsform (Organisational form)}\label{bylaw:3}
\textbf{\S 3.1 Uavhengighet (Independence).}\label{bylaw:3.1} Foreningen er politisk og religiøst uavhengig.

\textbf{\S 3.2 Juridisk status (Legal status).}\label{bylaw:3.2} Foreningen er en frittstående juridisk person med medlemmer, og er selveiende. Verken medlemmer eller andre har krav på foreningens formue eller eiendeler, og de er heller ikke ansvarlige for foreningens gjeld eller andre forpliktelser.

\textbf{\S 3.3 Driftsspråk (Operating language).}\label{bylaw:3.3}
\begin{enumerate}
  \item Foreningens driftsspråk er engelsk.
  \item Alle tekster foreningen produserer på norsk, skal oversettes til engelsk.
  \item Tekster på engelsk skal oversettes til norsk dersom det blir etterspurt.
\end{enumerate}

\vsubsection{\S 4 Definisjoner og ordforklaringer (Definitions)}\label{bylaw:4}
\textbf{\S 4.1 Virkeområde (Scope).}\label{bylaw:4.1} Begrepene i denne paragrafen gjelder gjennom hele vedtektene. Begrepene skal brukes konsekvent i vedtak, protokoller og annen organisatorisk dokumentasjon.

\textbf{\S 4.2 Organer (Organs).}\label{bylaw:4.2}
\begin{enumerate}
  \item Generalforsamling (GF) (General Assembly): foreningens øverste organ, jf. \hyperref[bylaw:8]{\S 8}.
  \item Styret (Board): foreningens øverste organ mellom generalforsamlingene, jf. \hyperref[bylaw:12]{\S 12}.
  \item Valg- og kontrollkomité (Election and Control Committee): en komité valgt av generalforsamlingen som gjennomfører valg til verv og overvåker at vedtektene følges, jf. \hyperref[bylaw:13]{\S 13}.
\end{enumerate}

\textbf{\S 4.3 Roller (Roles).}\label{bylaw:4.3}
\begin{enumerate}
  \item CEO (Chief Executive Officer): styreleder.
  \item CTO (Chief Technical Officer): teknisk ansvarlig og nestleder.
  \item CFO (Chief Financial Officer): forretnings- og økonomiansvarlig.
  \item CMO (Chief Marketing Officer): markedsføringsansvarlig.
  \item CPO (Chief People Officer): personal- og medlemsansvarlig.
  \item Prosjektleder (Project Lead): et medlem med overordnet ansvar for gjennomføringen av et prosjekt.
  \item Teamleder (Team Lead): et medlem med ansvar for et team, det vil si en faggruppe som går på tvers av flere prosjekter, jf. \hyperref[bylaw:12.9]{\S 12.9}.
  \item Nærmeste leder (Nearest Leader): den rollen et medlem rapporterer til i foreningens styringsstruktur. For et ordinært medlem er nærmeste leder teamlederen. Teamlederen rapporterer til prosjektlederen for det aktuelle prosjektet, og prosjektlederen til ansvarlig styremedlem, jf. \hyperref[bylaw:16.2]{\S 16.2}.
  \item Ansvarsverv (Responsibility Role): et tildelt og avgrenset lederansvar utover ordinært medlemskap, herunder prosjektleder og teamleder.
\end{enumerate}

\textbf{\S 4.4 Medlemskap (Membership).}\label{bylaw:4.4}
\begin{enumerate}
  \item Aktivt medlem (Active Member): synonymt med «fullt medlem». Begrepet aktivt medlem brukes i disse vedtektene, jf. \hyperref[bylaw:5]{\S 5}.
  \item Mentor (Mentor): medlemsnivå for personer som står på foreningens mentorliste, jf. \hyperref[bylaw:5.4]{\S 5.4}.
  \item Alumni (Alumni): medlemsnivå for personer hvis medlemskontrakt er utløpt uten fornyelse, jf. \hyperref[bylaw:5.4]{\S 5.4}.
  \item Aktiv periode (Active Period): NMBUs vår- og høstparallell, fra semesterstart til begynnelsen av eksamensperioden.
\end{enumerate}

\textbf{\S 4.5 Generalforsamling og valg (General Assembly and elections).}\label{bylaw:4.5}
\begin{enumerate}
  \item Valggeneralforsamling (Elective General Assembly): generalforsamlingen i vårsemesteret, der tillitsverv velges for kommende vervperiode, jf. \hyperref[bylaw:8.2]{\S 8.2}.
  \item Suppleringsgeneralforsamling (Supplementary General Assembly): generalforsamlingen i høstsemesteret, der det gjennomføres suppleringsvalg til ledige verv og valg til valg- og kontrollkomiteen, jf. \hyperref[bylaw:8.2]{\S 8.2}.
  \item Tillitsverv (Office): et verv som velges av generalforsamlingen for en vervperiode, herunder styreverv og verv i valg- og kontrollkomiteen, jf. \hyperref[bylaw:10]{\S 10}.
  \item Forslagsrett (Right to Propose): retten til å fremme saker og endringsforslag på generalforsamlingen, jf. \hyperref[bylaw:5.5]{\S 5.5}.
  \item Endringsforslag (Amendment): et forslag om å endre en sak på sakslisten, jf. \hyperref[bylaw:9.5]{\S 9.5}.
  \item Forhåndsbenking (Advance Nomination): benking av en kandidat før generalforsamlingen, jf. \hyperref[bylaw:11.3]{\S 11.3}.
  \item Benkeforslag (Floor Nomination): benking av en kandidat fra salen under generalforsamlingen, jf. \hyperref[bylaw:11.3]{\S 11.3}.
  \item Benketale (Floor Speech): en tale der den som har fremmet et benkeforslag, begrunner forslaget.
  \item Forhåndsstemme (Advance Vote): en stemme i et valg, avgitt før stemmegivningen i det aktuelle valget åpner på generalforsamlingen, jf. \hyperref[bylaw:11.4]{\S 11.4}.
\end{enumerate}

\textbf{\S 4.6 Møtefunksjonærer og møteregler (Meeting officers and meeting rules).}\label{bylaw:4.6}
\begin{enumerate}
  \item Ordstyrer (Meeting Chair): leder generalforsamlingen, jf. \hyperref[bylaw:9.2]{\S 9.2}.
  \item Referent (Secretary): fører protokollen for generalforsamlingen, jf. \hyperref[bylaw:9.8]{\S 9.8}.
  \item Tellekorps (Vote Counters): de medlemmene som teller stemmer på generalforsamlingen, jf. \hyperref[bylaw:9.3]{\S 9.3}.
  \item Protokollunderskriver (Protocol Signer): et medlem som kontrollerer og signerer protokollen, jf. \hyperref[bylaw:9.9]{\S 9.9}.
  \item Taleliste (Speakers' List): oversikten over hvem som har bedt om ordet, og i hvilken rekkefølge.
  \item Innlegg (Speech): en tale i diskusjonen av en sak, holdt i talelistens rekkefølge.
  \item Replikk (Reply): et kort svar på det umiddelbart foregående innlegget, jf. \hyperref[bylaw:9.4]{\S 9.4}.
  \item Innlegg til forretningsorden (Point of Order): et innlegg om møtets gjennomføring, som har forrang foran talelisten.
  \item Strek på talelisten (Closing the Speakers' List): et vedtak om at ingen nye talere kan tegne seg til saken.
  \item Voteringsorden (Voting Order): rekkefølgen forslagene settes under votering i, jf. \hyperref[bylaw:9.6]{\S 9.6}.
  \item Protokolltilførsel (Protocol Note): en kort merknad et medlem krever inntatt i protokollen, jf. \hyperref[bylaw:9.8]{\S 9.8}.
\end{enumerate}

\textbf{\S 4.7 Myndighet og fullmakt (Authority and mandate).}\label{bylaw:4.7}
\begin{enumerate}
  \item Myndighet (Authority): den beslutningsmyndigheten et organ eller en rolle har direkte etter disse vedtektene.
  \item Fullmakt (Mandate): retten til å treffe beslutninger og binde foreningen på CEOs vegne, innenfor et tildelt ansvarsområde, jf. \hyperref[bylaw:14]{\S 14}.
  \item Delegasjon (Delegation): overføring av fullmakt fra CEO til et annet medlem, jf. \hyperref[bylaw:14.2]{\S 14.2}.
  \item Signaturrett (Signature Right): retten til å signere bindende avtaler på foreningens vegne, jf. \hyperref[bylaw:14.7]{\S 14.7}.
  \item Føringsorden (Standing Order): regler vedtatt av styret som et tillegg til vedtektene, jf. \hyperref[bylaw:17.4]{\S 17.4}.
\end{enumerate}

\textbf{\S 4.8 Økonomi og prosjekter (Finance and projects).}\label{bylaw:4.8}
\begin{enumerate}
  \item Frie midler (Free Funds): midler foreningen disponerer uten øremerking, jf. \hyperref[bylaw:15.6]{\S 15.6}.
  \item NASA Project Life Cycle: rammeverket for faseinndeling og portgjennomganger (gate reviews) i prosjekter, i en modifisert versjon fastsatt i foreningens styringsdokumenter, jf. \hyperref[bylaw:16.3]{\S 16.3}.
  \item MCR (Mission Concept Review): konseptgjennomgangen som kreves for å starte et prosjekt, jf. \hyperref[bylaw:16.1]{\S 16.1}.
  \item SRR (System Requirements Review): kravgjennomgangen som kreves for å starte et prosjekt, jf. \hyperref[bylaw:16.1]{\S 16.1}.
\end{enumerate}

\vsection{Del II: Medlemskap (Part II: Membership)}

\vsubsection{\S 5 Medlemskap (Membership)}\label{bylaw:5}
\textbf{\S 5.1 Medlemsnivåer (Membership tiers).}\label{bylaw:5.1} Eclipse NMBU opererer med seks medlemsnivåer. Nivåene er hierarkiske, slik at høyere nivåer har alle rettigheter og plikter som tilhører nivåene under:
\begin{enumerate}
  \item \textbf{Styremedlem (Board Member):} har det juridiske og administrative ansvaret for foreningen.
  \item \textbf{Kjernemedlem (Core Member):} et medlem som er tildelt et ansvarsverv av styret eller et styremedlem, jf. \hyperref[bylaw:4.3]{\S 4.3} nr.\ 8, herunder prosjektleder, teamleder eller et annet avgrenset lederansvar.
  \item \textbf{Aktivt medlem (Active Member):} et medlem som er tatt opp etter \S5.2.
  \item \textbf{Mentor (Mentor):} en person som står på foreningens mentorliste, jf. \hyperref[bylaw:5.4]{\S 5.4}.
  \item \textbf{Inaktivt medlem (Inactive Member):} et medlem som er registrert i foreningen, men som ikke er aktivt medlem.
  \item \textbf{Alumni (Alumni):} et medlem med alumnistatus, jf. \hyperref[bylaw:5.4]{\S 5.4}.
\end{enumerate}
CPO avgjør om et medlem er aktivt eller inaktivt.

\textbf{\S 5.2 Opptak og medlemskontrakt (Admission and Membership Contract).}\label{bylaw:5.2}
\begin{enumerate}
  \item Alle medlemmer skal signere en medlemskontrakt som bekrefter at de har lest vedtektene, og at de forstår og aksepterer risikoen knyttet til romfartsvirksomhet.
  \item For å bli aktivt medlem må medlemmet ha gjennomført et intervju og signert medlemskontrakten.
  \item CPO er ansvarlig for opptaksprosessen. Intervjuet gjennomføres av to styremedlemmer, som avgjør opptaket.
\end{enumerate}

\textbf{\S 5.3 Medlemsperiode (Membership period).}\label{bylaw:5.3} Medlemsperioden følger vervperioden i \S10.1 og løper fra 1. august til 31. juli. En person kan tre inn når som helst i perioden, men medlemskapet opphører 31. juli. Mentorer og alumni er unntatt fra dette; deres medlemskap varer til vedkommende selv trekker seg, eller til CPO fjerner dem, jf. \hyperref[bylaw:5.4]{\S 5.4}.

\textbf{\S 5.4 Mentorer og alumni (Mentors and Alumni).}\label{bylaw:5.4}
\begin{enumerate}
  \item Mentor er et eget medlemsnivå, som også kan innehas av personer utenfor foreningen. Personer kan bli spurt om å bli mentor, og de som samtykker, føres opp på foreningens mentorliste. CPO er ansvarlig for mentorlisten.
  \item Et medlem hvis medlemskontrakt utløper uten å bli fornyet, får automatisk alumnistatus.
\end{enumerate}

\textbf{\S 5.5 Rettigheter på generalforsamlingen (Rights at the General Assembly).}\label{bylaw:5.5}
\begin{enumerate}
  \item Styremedlemmer, kjernemedlemmer og aktive medlemmer har stemmerett og er valgbare.
  \item Stemmerett forutsetter at medlemmet er tatt opp som aktivt medlem, jf. \S5.2, minst én måned før generalforsamlingen.
  \item Medlemmer med stemmerett, og mentorer, har forslagsrett.
  \item Alle medlemmer, også inaktive medlemmer, mentorer og alumni, har møte- og talerett.
  \item Inaktive medlemmer, mentorer og alumni har verken stemmerett eller valgbarhet.
  \item Generalforsamlingen kan med simpelt flertall innvilge møte- og talerett til personer som ikke er medlemmer.
\end{enumerate}

\textbf{\S 5.6 Plikter (Duties).}\label{bylaw:5.6}
\begin{enumerate}
  \item \textbf{Følgeplikt:} medlemmene er bundet av vedtak truffet av generalforsamlingen og styret.
  \item \textbf{Leveringsplikt:} et medlem som har påtatt seg en oppgave, er forpliktet til å levere den innenfor den fastsatte tidsrammen.
  \item \textbf{Varslingsplikt:} dersom et medlem ikke kan gjennomføre en oppgave det har påtatt seg, skal medlemmet uten ugrunnet opphold varsle sin nærmeste leder, jf. \hyperref[bylaw:4.3]{\S 4.3} nr.\ 7. Nærmeste leder er ansvarlig for at oppgaven blir fordelt på nytt.
  \item \textbf{Fravær:} dersom et medlem er inaktivt uten gyldig grunn i en sammenhengende periode på fire uker i en aktiv periode, kan styret vurdere å avslutte medlemskapet. Årsaken til fraværet skal avklares med medlemmets nærmeste leder.
\end{enumerate}

\textbf{\S 5.7 Styrets myndighet i medlemssaker (Board authority in membership matters).}\label{bylaw:5.7} Utenfor generalforsamlingen kan styret ved flertallsvedtak:
\begin{enumerate}
  \item Endre et medlems nivå på grunnlag av brudd på pliktene, likevel slik at vurderingen av om et medlem er aktivt eller inaktivt er unntatt, jf. \hyperref[bylaw:5.1]{\S 5.1}.
  \item Innvilge eller avslå søknader om verv og medlemskap, og endre hvem som innehar et verv.
  \item Fjerne medlemmer fra foreningens medlemsliste dersom \S5.6 nr.\ 4 brytes.
\end{enumerate}

\vsubsection{\S 6 Varsling og reaksjoner (Whistleblowing and reactions)}\label{bylaw:6}
\textbf{\S 6.1 Definisjon (Definition).}\label{bylaw:6.1} Varsling (whistleblowing) er når et medlem, et styremedlem eller en person i kontakt med foreningen melder om et kritikkverdig forhold knyttet til foreningens virksomhet.

\textbf{\S 6.2 Kritikkverdige forhold (Reprehensible matters).}\label{bylaw:6.2} Kritikkverdige forhold omfatter blant annet:
\begin{enumerate}
  \item Forhold som gjelder helse, miljø og sikkerhet.
  \item Personkonflikter.
  \item Forstyrrende oppførsel.
  \item Lovbrudd, herunder tyveri, økonomisk utroskap og underslag.
  \item Seksuell trakassering.
  \item Overgrep.
\end{enumerate}

\textbf{\S 6.3 Saksbehandling (Case handling).}\label{bylaw:6.3}
\begin{enumerate}
  \item Et varsel rettes til CPO, eller til CEO dersom varselet gjelder CPO.
  \item Ved alvorlige anklager om lovbrudd må saken anmeldes til politiet før styret kan behandle varselet.
  \item I saker som er rettet mot en eller flere personer, skal alle involverte gis mulighet til å redegjøre for sin versjon av situasjonen.
  \item Lovbrudd og seksuell trakassering aksepteres ikke, og styret kan anmelde forhold til politiet på grunnlag av egne observasjoner.
\end{enumerate}

\textbf{\S 6.4 Habilitet (Impartiality).}\label{bylaw:6.4} Ved behandling av varsler skal styrets habilitet vurderes. Et styremedlem som ikke er habilt, deltar ikke i behandlingen. Dersom styret som helhet ikke er habilt, skal det oppnevnes et uavhengig organ til å behandle saken.

\textbf{\S 6.5 Reaksjoner (Reactions).}\label{bylaw:6.5} Vedtak om en reaksjon krever 2/3 flertall i styret, jf. \hyperref[bylaw:12.7]{\S 12.7} nr.\ 5. Følgende reaksjoner kan besluttes:
\begin{enumerate}
  \item Endringer i foreningens rutiner.
  \item En advarsel.
  \item Suspensjon eller utestengelse fra verv og foreningens aktiviteter for en periode.
  \item Permanent utestengelse fra verv og foreningens aktiviteter.
\end{enumerate}

\vsubsection{\S 7 Personvern (Privacy)}\label{bylaw:7}
\textbf{\S 7.1 Behandling av personopplysninger (Processing of personal data).}\label{bylaw:7.1} Foreningen behandler bare personopplysninger som er nødvendige og relevante for å administrere medlemskap. Behandlingen skal følge personopplysningsloven og Datatilsynets retningslinjer.

\textbf{\S 7.2 Oppbevaring ved opphør av medlemskap (Retention on end of membership).}\label{bylaw:7.2} Ved opphør av medlemskap oppbevarer foreningen medlemmets personopplysninger til det fremsettes en anmodning om sletting etter \S7.3.

\textbf{\S 7.3 Anmodning om sletting (Request for deletion).}\label{bylaw:7.3}
\begin{enumerate}
  \item Et medlem eller tidligere medlem kan når som helst kreve at egne personopplysninger slettes.
  \item Anmodningen skal fremsettes skriftlig til CPO.
  \item Sletting skal gjennomføres uten ugrunnet opphold, og medlemmet skal motta bekreftelse på at slettingen er gjennomført.
\end{enumerate}

\textbf{\S 7.4 Informasjonsplikt (Duty to inform).}\label{bylaw:7.4} Foreningen skal ved opptak og ved opphør av medlemskap informere medlemmet om retten til å kreve sletting etter \S7.3.


\vsection{Del III: Organer (Part III: Organs)}

\vsubsection{\S 8 Generalforsamling (General Assembly)}\label{bylaw:8}
\textbf{\S 8.1 Øverste organ (Supreme organ).}\label{bylaw:8.1} Generalforsamlingen er foreningens øverste organ.

\textbf{\S 8.2 Hyppighet og typer (Frequency and types).}\label{bylaw:8.2}
\begin{enumerate}
  \item Generalforsamling skal avholdes minst én gang per semester, innen utgangen av november i høstsemesteret og innen utgangen av april i vårsemesteret.
  \item Generalforsamlingen i vårsemesteret er valggeneralforsamlingen; generalforsamlingen i høstsemesteret er suppleringsgeneralforsamlingen.
\end{enumerate}

\textbf{\S 8.3 Innkalling og frister (Notice and deadlines).}\label{bylaw:8.3}
\begin{enumerate}
  \item Innkalling skal være medlemmene i hende minst én måned før generalforsamlingen.
  \item Saker som skal behandles av generalforsamlingen, herunder forslag til vedtektsendringer, skal være styret i hende minst to uker før generalforsamlingen.
  \item Sakspapirer skal være tilgjengelige for medlemmene minst én uke før generalforsamlingen.
\end{enumerate}

\textbf{\S 8.4 Vedtaksførhet og flertallskrav (Quorum and majority requirements).}\label{bylaw:8.4}
\begin{enumerate}
  \item Generalforsamlingen er vedtaksfør når den er lovlig innkalt, uavhengig av antall møtende.
  \item Generalforsamlingen treffer vedtak med simpelt flertall med mindre disse vedtektene bestemmer noe annet. Flertallskrav for vedtektsendringer og oppløsning følger \S\S17.1 og 18.1.
  \item Ingen medlemmer har mer enn én stemme, og det kan ikke stemmes ved fullmektig. Forhåndsstemming i valg reguleres av \S11.4.
\end{enumerate}

\textbf{\S 8.5 Ekstraordinær generalforsamling (Extraordinary General Assembly).}\label{bylaw:8.5} Ekstraordinær generalforsamling kan innkalles dersom minst 1/3 av medlemmene krever det, eller dersom styret finner det nødvendig. \S\S8.3 og 8.4 gjelder tilsvarende. På en ekstraordinær generalforsamling behandles bare de sakene som er angitt i innkallingen.

\textbf{\S 8.6 Saksliste (Agenda).}\label{bylaw:8.6} Sakslistene for valg- og suppleringsgeneralforsamlingen følger rekkefølgen nedenfor.

\emph{Valggeneralforsamling (vårsemesteret):}
\begin{enumerate}
  \item Konstituering: (1) valg av ordstyrer; (2) valg av referent; (3) valg av tellekorps; (4) valg av protokollunderskrivere; (5) godkjenning av innkalling; (6) godkjenning av saksliste.
  \item Presentasjon og diskusjon av saker: (1) semesterrapport; (2) forrige semesters regnskap; (3) kommende semesters budsjett; (4) innsendte saker.
  \item Votering over saker.
  \item Valg av CEO, CTO, CFO, CMO og CPO for kommende vervperiode, jf. \hyperref[bylaw:11]{\S 11}.
  \item Eventuelt.
\end{enumerate}

\emph{Suppleringsgeneralforsamling (høstsemesteret):}
\begin{enumerate}
  \item Konstituering: (1) valg av ordstyrer; (2) valg av referent; (3) valg av tellekorps; (4) valg av protokollunderskrivere; (5) godkjenning av innkalling; (6) godkjenning av saksliste.
  \item Presentasjon og diskusjon av saker: (1) semesterrapport; (2) forrige semesters regnskap; (3) kommende semesters budsjett; (4) innsendte saker.
  \item Votering over saker.
  \item Suppleringsvalg, jf. \hyperref[bylaw:11]{\S 11}: (1) suppleringsvalg til ledige verv for resten av vervperioden; (2) valg av leder og to medlemmer til valg- og kontrollkomiteen.
  \item Eventuelt.
\end{enumerate}

Hver sak under punkt 2 presenteres og diskuteres samlet, og sakene voteres over under punkt 3. Saksbehandlingen følger \S9.

\vsubsection{\S 9 Møteregler for generalforsamlingen (Meeting rules for the General Assembly)}\label{bylaw:9}
\textbf{\S 9.1 Åpning og konstituering (Opening and constitution).}\label{bylaw:9.1}
\begin{enumerate}
  \item Generalforsamlingen åpnes av CEO, eller av CTO etter \S12.3. Den som åpner møtet, leder valget av ordstyrer.
  \item Inntil tellekorpset er valgt, telles stemmene av den som leder møtet.
  \item Når ordstyreren er valgt, overtar ordstyreren møteledelsen og gjennomfører de gjenstående punktene i konstitueringen.
  \item Ved godkjenning av innkallingen vurderer generalforsamlingen om innkallingen og sakspapirene er i samsvar med \S8.3. Dersom innkallingen ikke godkjennes, kan generalforsamlingen ikke treffe vedtak.
\end{enumerate}

\textbf{\S 9.2 Ordstyrer (Meeting Chair).}\label{bylaw:9.2}
\begin{enumerate}
  \item Ordstyreren leder generalforsamlingen i samsvar med vedtektene og den godkjente sakslisten, og skal sikre en ordnet og nøytral gjennomføring av møtet.
  \item Ordstyreren fører talelisten, gir ordet, og kan ta ordet fra en taler dersom taletiden overskrides eller innlegget ikke angår saken.
  \item Ordstyreren kan fastsette taletid for innlegg og replikker.
  \item Ordstyreren avgjør spørsmål om møtets gjennomføring. Avgjørelsen kan bringes inn for generalforsamlingen, som avgjør spørsmålet med simpelt flertall.
  \item Ordstyreren deltar ikke i debatten. Dersom ordstyreren ønsker å delta i debatten i en sak, overlates møteledelsen for den saken til en person valgt av generalforsamlingen.
  \item Ordstyreren bør ikke være et styremedlem.
\end{enumerate}

\textbf{\S 9.3 Tellekorps (Vote Counters).}\label{bylaw:9.3}
\begin{enumerate}
  \item Tellekorpset består av to medlemmer, som skal være nøytrale i de sakene de teller.
  \item Tellekorpset registrerer antallet stemmeberettigede som er til stede ved konstitueringen, og holder antallet oppdatert gjennom møtet.
  \item Tellekorpset deler ut og samler inn stemmesedler, kontrollerer at bare stemmeberettigede avgir stemme, teller stemmene og melder resultatet til ordstyreren.
  \item Tellekorpset teller stemmer i saker, og i valg som ikke gjennomføres av valg- og kontrollkomiteen, jf. \hyperref[bylaw:11.5]{\S 11.5}.
  \item Stemmesedlene oppbevares av tellekorpset til protokollen er signert, og makuleres deretter.
\end{enumerate}

\textbf{\S 9.4 Presentasjon og diskusjon av saker (Presentation and discussion of Cases).}\label{bylaw:9.4}
\begin{enumerate}
  \item Hver sak presenteres av forslagsstilleren eller styret, og diskuteres før neste sak presenteres.
  \item Under diskusjonen gis ordet i den rekkefølgen det er bedt om.
  \item Etter et innlegg kan det gis replikk. Den som holdt innlegget, har rett til én replikk.
  \item Innlegg til forretningsorden har forrang foran talelisten.
  \item Ordstyreren, eller enhver stemmeberettiget, kan foreslå strek på talelisten. Generalforsamlingen avgjør forslaget med simpelt flertall.
\end{enumerate}

\textbf{\S 9.5 Endringsforslag (Amendments).}\label{bylaw:9.5} Et endringsforslag til en sak på sakslisten må leveres skriftlig før voteringen over saker begynner.

\textbf{\S 9.6 Stemming over saker (Voting on Cases).}\label{bylaw:9.6}
\begin{enumerate}
  \item Før voteringen legger ordstyreren frem en voteringsorden, som godkjennes av generalforsamlingen.
  \item Endringsforslag voteres over før det opprinnelige forslaget. Dersom det foreligger flere endringsforslag til samme punkt, voteres det over det mest omfattende forslaget først.
  \item Det voteres ved håndsopprekning, med mindre minst én stemmeberettiget krever skriftlig og hemmelig votering.
  \item Når voteringen er i gang, gis ordet bare til spørsmål om voteringen selv.
  \item Blanke stemmer regnes som ikke avgitt.
  \item Ved stemmelikhet anses forslaget som ikke vedtatt.
\end{enumerate}

\textbf{\S 9.7 Eventuelt (Other business).}\label{bylaw:9.7}
\begin{enumerate}
  \item Saker til eventuelt må leveres skriftlig før voteringen over saker begynner.
  \item Under eventuelt behandles bare saker som er levert etter nr.\ 1.
  \item Det kan ikke treffes vedtak om vedtektsendringer, valg, økonomiske disposisjoner eller oppløsning under eventuelt.
\end{enumerate}

\textbf{\S 9.8 Referent og protokoll (Secretary and protocol).}\label{bylaw:9.8}
\begin{enumerate}
  \item Referenten fører protokollen for generalforsamlingen.
  \item Protokollen skal minst inneholde: (1) tid og sted for møtet; (2) antall stemmeberettigede til stede; (3) hvem som ble valgt under konstitueringen; (4) alle forslag som ble fremmet, med forslagsstiller; (5) vedtak med stemmetall; (6) resultatet av alle valg; (7) protokolltilførsler.
  \item Enhver stemmeberettiget kan kreve en kort protokolltilførsel inntatt i protokollen.
  \item Protokollen skal signeres av protokollunderskriverne og gjøres tilgjengelig for medlemmene senest to uker etter generalforsamlingen.
\end{enumerate}

\textbf{\S 9.9 Protokollunderskrivere (Protocol Signers).}\label{bylaw:9.9}
\begin{enumerate}
  \item Generalforsamlingen velger to protokollunderskrivere blant de møtende medlemmene.
  \item Protokollunderskriverne kontrollerer at protokollen gir en korrekt fremstilling av generalforsamlingen, og signerer den. De kan kreve rettelser før de signerer.
\end{enumerate}

\vsubsection{\S 10 Tillitsverv (Offices)}\label{bylaw:10}
\textbf{\S 10.1 Vervperiode (Term of office).}\label{bylaw:10.1} Alle verv i foreningen, herunder alle styreverv, løper fra 1. august til 31. juli året etter. Valg- og kontrollkomiteen er unntatt fra dette, jf. \hyperref[bylaw:13.1]{\S 13.1}.

\textbf{\S 10.2 Overlapp og erfaringsoverføring (Overlap and experience transfer).}\label{bylaw:10.2} Perioden fra valggeneralforsamlingen til 1. august er en overlappsperiode der påtroppende og avtroppende vervinnehaver fungerer samtidig. I denne perioden har den påtroppende vervinnehaveren beslutningsmyndighet og stemmerett, mens den avtroppende vervinnehaveren har en rådgivende funksjon.

Den avtroppende vervinnehaveren skal i overlappsperioden sørge for full erfaringsoverføring, herunder overlevering av dokumentasjon, tilganger, løpende saker og kontaktpunkter, og skal være tilgjengelig for styret og sin etterfølger.

\textbf{\S 10.3 Ledige verv (Vacant offices).}\label{bylaw:10.3}
\begin{enumerate}
  \item Dersom en vervinnehaver trekker seg før vervperioden er utløpt, kan styret med simpelt flertall konstituere et medlem i vervet til det blir fylt på en ordinær generalforsamling.
  \item Dersom et verv ikke blir fylt på generalforsamlingen, kan styret konstituere et medlem i vervet til neste generalforsamling.
  \item Nr.\ 1 og 2 gjelder tilsvarende for valg- og kontrollkomiteen. Dersom ingen velges til komiteen, kan styret oppnevne en midlertidig komité til neste generalforsamling.
\end{enumerate}

\vsubsection{\S 11 Valg (Elections)}\label{bylaw:11}
\textbf{\S 11.1 Virkeområde (Scope).}\label{bylaw:11.1} Denne paragrafen gjelder for alle valg som holdes på generalforsamlingen, herunder valg av ordstyrer, referent, tellekorps, protokollunderskrivere, styremedlemmer og medlemmer av valg- og kontrollkomiteen.

\textbf{\S 11.2 Kandidater (Candidates).}\label{bylaw:11.2}
\begin{enumerate}
  \item Stemmerett og valgbarhet følger \S5.5.
  \item En kandidat må ha samtykket til kandidaturet. Kandidater som ikke kan møte, kan stille med et skriftlig kandidatur.
\end{enumerate}

\textbf{\S 11.3 Forhåndsbenking og benkeforslag (Advance Nomination and Floor Nomination).}\label{bylaw:11.3}
\begin{enumerate}
  \item Før generalforsamlingen sendes det ut en oppfordring til forhåndsbenking. Valg- og kontrollkomiteen spør de foreslåtte om de ønsker å stille. De som ønsker å stille, føres opp som kandidater på sakslisten. Fristen for forhåndsbenking er den samme som for innsending av saker, jf. \hyperref[bylaw:8.3]{\S 8.3}.
  \item Under generalforsamlingen kan benkeforslag bare fremmes under presentasjon og diskusjon av saker, jf. \hyperref[bylaw:8.6]{\S 8.6}.
  \item Den som fremmer et benkeforslag, har rett til å holde en benketale.
\end{enumerate}

\textbf{\S 11.4 Forhåndsstemming (Advance Voting).}\label{bylaw:11.4}
\begin{enumerate}
  \item Aktive medlemmer med stemmerett kan avgi forhåndsstemme i valg til verv. Forhåndsstemming gjelder ikke andre saker.
  \item Forhåndsstemmer avgis digitalt gjennom en løsning som sikrer at bare stemmeberettigede kan stemme, at hvert medlem avgir én stemme, og at stemmen er hemmelig.
  \item Fristen for forhåndsstemming er når stemmegivningen i det aktuelle valget åpner på generalforsamlingen.
  \item Et medlem som har avgitt forhåndsstemme, kan ikke stemme på nytt i samme valg.
  \item Valg- og kontrollkomiteen administrerer forhåndsstemmene og teller dem sammen med stemmene avgitt på generalforsamlingen.
  \item En forhåndsstemme gjelder for den kandidaten den er avgitt for, også dersom nye kandidater kommer til ved benkeforslag på generalforsamlingen.
  \item I en andre avstemningsrunde teller forhåndsstemmen dersom kandidaten den er avgitt for, deltar i runden. Ellers teller den ikke.
\end{enumerate}

\textbf{\S 11.5 Gjennomføring (Conduct).}\label{bylaw:11.5}
\begin{enumerate}
  \item Ved valg til verv leder valg- og kontrollkomiteen valget og presenterer vervet og kandidatene. Andre valg ledes av ordstyreren, med unntak av valget av ordstyrer selv, jf. \hyperref[bylaw:9.1]{\S 9.1}.
  \item Kandidatene gis mulighet til å presentere seg, og salen mulighet til å stille spørsmål.
  \item Valget gjennomføres skriftlig når det er mer enn én kandidat til et verv, eller når minst én stemmeberettiget krever det. Ellers kan valget gjennomføres ved akklamasjon.
  \item Ved valg til verv teller valg- og kontrollkomiteen stemmene og melder resultatet til ordstyreren. For andre valg telles stemmene etter \S\S9.1 og 9.3. Resultatet føres i protokollen.
\end{enumerate}

\textbf{\S 11.6 Avstemningsrunder (Voting rounds).}\label{bylaw:11.6}
\begin{enumerate}
  \item Den første avstemningsrunden gjennomføres med blank stemme som et alternativ.
  \item Den kandidaten som får flest stemmer, er valgt.
  \item Dersom alternativet blank stemme får flest stemmer, gjennomføres en andre avstemningsrunde mellom de to kandidatene med flest stemmer, uten alternativet blank stemme.
  \item Dersom to eller flere kandidater står likt med flest stemmer, gjennomføres en andre avstemningsrunde mellom disse kandidatene, uten alternativet blank stemme.
  \item Dersom ingen kandidat blir valgt i den andre avstemningsrunden, gjentas den til en kandidat er valgt.
\end{enumerate}

\vsubsection{\S 12 Styret (The Board)}\label{bylaw:12}
\textbf{\S 12.1 Sammensetning (Composition).}\label{bylaw:12.1} Styret består av fire medlemmer:
\begin{enumerate}
  \item CEO: Chief Executive Officer
  \item CTO: Chief Technical Officer
  \item CFO: Chief Financial Officer
  \item CMO: Chief Marketing Officer
  \item CPO: Chief People Officer
\end{enumerate}
Alle styremedlemmer velges av generalforsamlingen, jf. \hyperref[bylaw:1]{\S 1}1, for vervperioden i \S10.1.

\textbf{\S 12.2 CEO.}\label{bylaw:12.2} CEO har den overordnede administrative ledelsen av foreningen og representerer foreningen utad. CEO har det endelige ansvaret for at foreningen opererer innenfor norsk lov og universitetets regelverk, og for strategisk planlegging, organisasjonskultur, HMS-protokoller og koordinering mellom de øvrige styrerollene. CEOs fullmakt og delegasjonsplikt følger \S14.

\textbf{\S 12.3 CTO.}\label{bylaw:12.3} CTO har det overordnede ansvaret for foreningens tekniske prosjekter, herunder systemintegrasjon, teknisk sikkerhet og teknisk fremdrift. CTO skal sikre at den tekniske utviklingen følger fastsatte tidslinjer og milepæler, jf. \hyperref[bylaw:1]{\S 1}6, og er ansvarlig for ressursforvaltning og logistikk knyttet til teknisk utstyr og verkstedsareal. CTO innkaller til tekniske statusmøter.

CTO er foreningens nestleder og overtar CEOs ansvar når CEO er utilgjengelig. Utover dette har CTO ikke større myndighet enn de øvrige styremedlemmene.

\textbf{\S 12.4 CFO.}\label{bylaw:12.4-cfo} CFO er ansvarlig for foreningens økonomi, jf. \hyperref[bylaw:15]{\S 15}. Dette omfatter økonomisk planlegging, bokføring, budsjettkontroll, revisjon, forsikring og finansiell risikostyring. CFO har videre ansvaret for sponsorrelasjoner, samarbeidsavtaler, merkevarebygging og ekstern kommunikasjon, herunder innhold på og vedlikehold av foreningens digitale plattformer.

\textbf{\S 12.4 CMO.}\label{bylaw:12.4-cmo} CMO er ansvarlig for ansvaret for sponsorrelasjoner, samarbeidsavtaler, merkevarebygging og ekstern kommunikasjon, herunder innhold på og vedlikehold av foreningens digitale og fysiske plattformer.

\textbf{\S 12.5 CPO.}\label{bylaw:12.5} CPO er ansvarlig for foreningens medlemmer og interne organisasjon, herunder rekruttering, opptaksprosessen, medlemsoppfølging og administrasjon av medlemsnivåene, jf. \hyperref[bylaw:5]{\S 5}. CPO har videre ansvaret for intern kommunikasjon, organisasjonskultur og oppfølging av HMS, i samarbeid med CEO, og for saksbehandling av varsler, jf. \hyperref[bylaw:6]{\S 6}.

\textbf{\S 12.6 Styremøter (Board meetings).}\label{bylaw:12.6} Styret er vedtaksført når minst tre av styrets medlemmer, inkludert konstituerte medlemmer, er til stede. Vedtak treffes med alminnelig flertall av de avgitte stemmene, med mindre disse vedtektene bestemmer noe annet. Alle styremøter skal føres i protokoll.

\textbf{\S 12.7 Stemmegivning i styret (Voting on the Board).}\label{bylaw:12.7}
\begin{enumerate}
  \item Hvert styremedlem har én stemme. CEOs stemme teller som 1,5 stemmer i alle styresaker.
  \item Når CTO har overtatt CEOs ansvar etter \S12.3, teller CTOs stemme som 1,5 stemmer. Den ekstra vektingen på 0,5 kan ikke tilfalle andre styremedlemmer.
  \item Et konstituert styremedlem har stemmerett, jf. \hyperref[bylaw:10.3]{\S 10.3}.
  \item Ved stemmelikhet anses forslaget som ikke vedtatt.
  \item Der disse vedtektene krever 2/3 flertall i styret, beregnes flertallet ut fra den samlede stemmevekten til de styremedlemmene som deltar i voteringen.
\end{enumerate}

\textbf{\S 12.8 Lojalitetsplikt (Duty of loyalty).}\label{bylaw:12.8} Styrets medlemmer, inkludert CEO, skal i styresaker sette foreningens samlede interesser foran andre hensyn.

\textbf{\S 12.9 Team (Teams).}\label{bylaw:12.9}
\begin{enumerate}
  \item Team opprettes av det styremedlemmet som er ansvarlig for fagområdet, og går på tvers av flere prosjekter.
  \item Styremedlemmet utpeker en teamleder for teamet.
\end{enumerate}

\vsubsection{\S 13 Valg- og kontrollkomiteen (The Election and Control Committee)}\label{bylaw:13}
\textbf{\S 13.1 Sammensetning og vervperiode (Composition and term).}\label{bylaw:13.1}
\begin{enumerate}
  \item Valg- og kontrollkomiteen består av tre medlemmer: en leder og to medlemmer.
  \item Lederen og medlemmene velges på suppleringsgeneralforsamlingen i høstsemesteret.
  \item Vervperioden løper fra 1. januar etter suppleringsgeneralforsamlingen til 31. desember samme år.
  \item Styremedlemmer kan ikke være medlemmer av komiteen.
\end{enumerate}

\textbf{\S 13.2 Oppgaver (Tasks).}\label{bylaw:13.2}
\begin{enumerate}
  \item Komiteen skal finne egnede kandidater til valg, administrere forhåndsstemmer og gjennomføre valg til verv, jf. \hyperref[bylaw:11]{\S 11}.
  \item Komiteen skal overvåke at vedtektene følges, og skal rapportere til generalforsamlingen.
\end{enumerate}

\textbf{\S 13.3 Taushetsplikt og valgbarhet (Duty of confidentiality and eligibility).}\label{bylaw:13.3}
\begin{enumerate}
  \item Komiteens medlemmer har taushetsplikt.
  \item Medlemmer av komiteen kan ikke stille til valg.
\end{enumerate}


\vsubsection{\S 14 Fullmakt, delegasjon og signaturrett (Mandate, delegation and Signature Right)}\label{bylaw:14}
\textbf{\S 14.1 Utgangspunkt (Starting point).}\label{bylaw:14.1} CEO har foreningens overordnede fullmakt innenfor de rammene generalforsamlingen og styret setter, og kan inngå bindende avtaler på foreningens vegne innenfor disse rammene, jf. \S\S14.7 og 14.8.

\textbf{\S 14.2 Delegasjonsplikt (Duty to delegate).}\label{bylaw:14.2} CEO skal delegere fullmakt til de øvrige styremedlemmene innenfor deres respektive ansvarsområder, jf. \S\S12.3--12.5. Delegasjonen skal legges frem for, og godkjennes av, generalforsamlingen. Myndigheten over foreningens økonomi tilligger CFO direkte, jf. \hyperref[bylaw:1]{\S 1}5.2, og omfattes ikke av delegasjonen. Delegasjon av signaturrett følger \S14.9.

\textbf{\S 14.3 Form og innhold (Form and content).}\label{bylaw:14.3} Delegasjonen skal være skriftlig og skal angi hvilke ansvarsområder, beslutningskategorier og eventuelle økonomiske grenser den omfatter. Delegasjonen gjelder for vervperioden, jf. \hyperref[bylaw:1]{\S 1}0.1, med mindre noe annet er bestemt, og skal føres i generalforsamlingens protokoll.

\textbf{\S 14.4 Ansvar (Responsibility).}\label{bylaw:14.4} Delegasjon fritar ikke CEO for det overordnede ansvaret for foreningens virksomhet.

\textbf{\S 14.5 Endring og tilbakekalling (Change and withdrawal).}\label{bylaw:14.5} Endring eller tilbakekalling av delegert fullmakt behandles av generalforsamlingen. Der forholdene ikke kan vente på en ordinær generalforsamling, kan styret med 2/3 flertall midlertidig suspendere en delegert fullmakt til neste generalforsamling.

\textbf{\S 14.6 Videredelegering (Onward delegation).}\label{bylaw:14.6} Et styremedlem kan videredelegere avgrensede deler av sin fullmakt til et aktivt medlem for bestemte oppgaver. Videredelegering skal være skriftlig og meldes til styret, og fritar ikke styremedlemmet for ansvar.

\textbf{\S 14.7 Signaturrett (Signature Right).}\label{bylaw:14.7} CEO har foreningens signaturrett.

\textbf{\S 14.8 Vilkår for signering (Conditions for signing).}\label{bylaw:14.8} Signering av en bindende avtale forutsetter et forutgående styrevedtak.

\textbf{\S 14.9 Delegasjon av signaturrett (Delegation of Signature Right).}\label{bylaw:14.9} CEO kan delegere signaturrett direkte til de øvrige styremedlemmene, uten behandling i generalforsamlingen. Delegasjonen skal dokumenteres skriftlig og fritar ikke CEO for ansvar, jf. \hyperref[bylaw:14.4]{\S 14.4}.

\textbf{\S 14.10 Ansvar for kontrakter (Responsibility for contracts).}\label{bylaw:14.10} Det sittende styret er ansvarlig for alle avtaler som er i kraft.

\textbf{\S 14.11 Digital oppbevaring (Digital storage).}\label{bylaw:14.11} Alle foreningens dokumenter, herunder avtaler og kontrakter som er signert på papir, skal også lagres digitalt. Et dokument som er signert på papir, skal skannes eller på annen måte gjengis i digital form uten ugrunnet opphold; den fysiske signaturen fritar ikke dokumentet fra dette kravet.

\vsection{Del IV: Drift (Part IV: Operations)}

\vsubsection{\S 15 Økonomi (Economy)}\label{bylaw:15}
\textbf{\S 15.1 Bruk av midler (Use of funds).}\label{bylaw:15.1} Foreningens midler skal bare brukes i samsvar med foreningens formål, jf. \hyperref[bylaw:2]{\S 2}.

\textbf{\S 15.2 Myndighet over økonomien (Authority over the economy).}\label{bylaw:15.2} CFO har full myndighet over foreningens økonomi. Denne myndigheten følger direkte av disse vedtektene og er ikke delegert fra CEO, jf. \hyperref[bylaw:14.2]{\S 14.2}.

\textbf{\S 15.3 Budsjett og regnskap (Budget and accounts).}\label{bylaw:15.3} Det skal til enhver tid finnes et budsjett og et regnskap for foreningens midler, med tilhørende bilag. Budsjettet og regnskapet gjennomgås og godkjennes av generalforsamlingen, jf. \hyperref[bylaw:8.6]{\S 8.6}.

\textbf{\S 15.4 Godkjenning av utbetalinger (Approval of payments).}\label{bylaw:15.4} Utbetalinger godkjennes av CFO.

\textbf{\S 15.5 Økonomiske investeringer (Financial investments).}\label{bylaw:15.5} Vedtak om økonomiske investeringer krever 2/3 flertall i styret, jf. \hyperref[bylaw:12.7]{\S 12.7} nr.\ 5.

\textbf{\S 15.6 Frie midler (Free Funds).}\label{bylaw:15.6}
\begin{enumerate}
  \item Midler foreningen tjener som ikke er øremerket ved mottak, skal plasseres på en egen konto for frie midler.
  \item Medlemmer kan søke om støtte fra frie midler til prosjekter. Søknaden rettes til CFO.
  \item CFO vurderer søknadene og avgjør hvilke som sendes videre til styret. Styret vedtar tildelingen.
  \item 75\% av midler fra en konkurranse tilordnes det aktuelle teamet med førsterett i ett år fra midlene er innvilget, etter CFOs og styrets skjønn med hensyn til hvordan midlene er tilgjengelige for bruk.
\end{enumerate}

\textbf{\S 15.7 Rapportering (Reporting).}\label{bylaw:15.7} CFO skal følge opp alle økonomiske forhold og holde styret løpende informert om foreningens økonomiske disposisjoner.

\vsubsection{\S 16 Prosjektgjennomføring (Project execution)}\label{bylaw:16}
\textbf{\S 16.1 Oppstart av prosjekt (Starting a project).}\label{bylaw:16.1} Alle større prosjekter skal godkjennes av styret før oppstart. For å starte et prosjekt skal det leveres en MCR og en SRR til styret. Styret vurderer prosjektet mot foreningens formål, ressurssituasjon og gjennomførbarhet, og godkjenner eller avslår det i det møtet der MCR og SRR presenteres.

\textbf{\S 16.2 Prosjektorganisering (Project organisation).}\label{bylaw:16.2} Styret utpeker en prosjektleder for hvert prosjekt. Prosjektlederen rapporterer til CTO i tekniske prosjekter, og ellers til det styremedlemmet som er ansvarlig for fagområdet.

\textbf{\S 16.3 Prosjektmodell (Project model).}\label{bylaw:16.3} Prosjekter gjennomføres etter en modifisert versjon av NASA Project Life Cycle. Faseinndeling, portgjennomganger og dokumentasjonskrav er fastsatt i foreningens styringsdokumenter.

\textbf{\S 16.4 Avvik (Deviations).}\label{bylaw:16.4} Vesentlige avvik fra en godkjent plan, herunder endringer i krav, budsjett, tidslinje eller sikkerhetsrelevante forutsetninger, skal legges frem for styret.

\textbf{\S 16.5 Avslutning og dokumentasjon (Closure and documentation).}\label{bylaw:16.5} Et prosjekt avsluttes med en sluttrapport til styret. Prosjektdokumentasjonen skal arkiveres slik at senere medlemmer kan bygge videre på arbeidet, jf. \hyperref[bylaw:2.5]{\S 2.5}.

\vsection{Del V: Styringsdokumenter og opphør (Part V: Governance Documents and Dissolution)}

\vsubsection{\S 17 Vedtektsendringer og styringsdokumenter (Bylaw Amendments and governance documents)}\label{bylaw:17}
\textbf{\S 17.1 Vedtektsendringer (Bylaw Amendments).}\label{bylaw:17.1} Vedtektsendringer kan bare skje når (1) minst 2/3 av foreningens medlemmer er til stede på generalforsamlingen på voteringstidspunktet, og (2) minst 2/3 av de møtende medlemmene stemmer for endringen. Et forslag til vedtektsendring må leveres innen fristen i \S8.3.

\textbf{\S 17.2 Ikrafttredelse (Entry into force).}\label{bylaw:17.2} Vedtektsendringer trer i kraft når protokollen fra generalforsamlingen er signert, med mindre generalforsamlingen vedtar noe annet.

\textbf{\S 17.3 Språkversjoner (Language versions).}\label{bylaw:17.3} Ved konflikt mellom språkversjoner av vedtektene er den versjonen generalforsamlingen har vedtatt, avgjørende, jf. \hyperref[bylaw:3.3]{\S 3.3}.

\textbf{\S 17.4 Føringsorden (Standing Orders).}\label{bylaw:17.4} Styret kan vedta en føringsorden som et tillegg til vedtektene. En føringsorden kan ikke være i strid med vedtektene, og faller bort ved styreskifte.

\vsubsection{\S 18 Opphør (Dissolution)}\label{bylaw:18}
\textbf{\S 18.1 Oppløsning (Dissolution).}\label{bylaw:18.1} Foreningen kan oppløses dersom 2/3 av de stemmeberettigede medlemmene på generalforsamlingen stemmer for oppløsning. En sak om oppløsning skal leveres innen fristen i \S8.3.

\textbf{\S 18.2 Disponering av eiendeler (Disposal of assets).}\label{bylaw:18.2} Ved oppløsning overføres foreningens midler og eiendeler til en annen egnet studentforening med samme formål, valgt av generalforsamlingen. Ingen medlemmer har krav på foreningens midler eller noen andel av dem, jf. \hyperref[bylaw:3.2]{\S 3.2}.

\end{document}
